\documentclass{article}

\usepackage{../handouts}

\title{CSCI 432 Handout 07: Decrementing Functions}
\author{Name: \rule{3in}{0.15mm} \\~\\
    Collaborators: \rule{3in}{0.15mm} }
\date{5 February 2026}



\begin{document}
\maketitle

\noindent
\section*{Definitions}

\begin{definition}[State Space, $\mathcal{S}$]
    A \emph{state} is a snapshot in the execution of an algorithm. It includes
    all states/values of all variables, stack traces, the number of times a loop
    has executed, etc. Think of this as everything that you can see uising a
    breakpoint when running your code through a debugger.  Note that some
    variables might be implicit (e.g., number of times a while loop has
    executed).

    The \emph{state space} of a program is the set of all states that are
    realized during the execution of the program. We call this a \emph{space} and not
    a \emph{set} as we can add edges between ``adjacent'' states to give the set
    more structure (namely, that of a digraph). Here, we just care about the set
    of all states though.
\end{definition}

\begin{definition}[Well-Ordered Set]
    A \emph{well-ordered set} is a partially ordered set (poset)
    such that every subset has a minimum or least
    element.  Mathematically, $(S, \leq)$ is a well-ordered set iff for all $S'
    \subset S$, $\exists! e \in S'$ such that for all $s \in S'$, $e \leq s'$.

    Our favorite well-ordered sets are $\N:=\{0,1,2, \ldots\}$ and
    $\Z_+ :=\{1,2,3, \ldots\}$.
\end{definition}

\section*{Decrementing Functions}

A decrementing function is used to prove that an algorithm (specifically, a loop
or a recurrence) terminates.  We use
them in two ways: (1) when doing back-of-hand calculations to ensure that our
algorithm will terminate before doing a proper runtime analysis; (2) when it is
difficult to prove exact runtime.

\begin{definition}[Decrementing Function]
    A decrementing function is a function $D \colon \mathcal{S} \to \N$ that
    satisfies the following properties:
    \begin{enumerate}[(1)]
        \item Each time a recursive call is made, the function strictly
            decreases between the top of the parent call to the top of the child
            call.
        \item Each time a loop is re-entered, the function is strictly less than
            the last time it entered that loop.
    \end{enumerate}
\end{definition}

If such a function exists, then our algorithm terminates!  Can you explain why
this is true?

\pagebreak
\section*{Practice}

Provide decrementing functions for the following algorithms:

\begin{enumerate}
    \item Suppose the algorithm has a single for loop, which appears as
        follows:
        \texttt{for i=10; i >0; i--}.
        What is the decrementing function?
        \practice
    \item Write a for loop that iterates through an array $A$ of length $n$
        starting at $A[1]$ and checking each index in order until $A[n]$, in
        order to find the maximum value in the array.  Then, provide the
        decrementing function for that loop.
        \practice
    \item Suppose there is a perfect infinite neighborhood, where the post
        office is at $(0,0)$ and for every $x, y \in \N$, there is a house at
        $(x,y)$. If you a person in this neighborhood can follow simple
        instructions (e.g., walk one block north/south/east/west), write a while
        loop that will give directions from a given house to the post office.
        Then, prove that your algorithm terminates.
        \practice
    \item Bubble sort. (Note: there are two loops here! On first attempt, just
        worry about the loops independently and come up with two functions.)

\end{enumerate}

\end{document}
